% Source created because the repository contains no TerraWatch Beamer/LaTeX source.
% It is intentionally a standalone, evidence-bounded update deck; the existing
% final.pdf at the repository root is a different RESQ flood-communications deck.
\documentclass[aspectratio=169,10pt]{beamer}
\usetheme{default}
\usecolortheme{dove}
\definecolor{canopy}{HTML}{0F766E}
\definecolor{ink}{HTML}{18312D}
\setbeamercolor{structure}{fg=canopy}
\setbeamercolor{normal text}{fg=ink}
\setbeamertemplate{navigation symbols}{}
\title{TerraWatch: Evidence-Aware Wildfire Workflows}
\subtitle{Research and dashboard integration update}
\author{TerraWatch project}
\date{October 2026}

\begin{document}
\begin{frame}\titlepage\end{frame}

\begin{frame}{Fire Image Authenticity Assessment}
  \begin{itemize}
    \item Existing assessor: CLIP ViT-B/16 QuickGELU with rank-8 LoRA; the shipped authenticity adapter uses block-8 CLS features. The reported source-shift results primarily support a final-layer variant, so they must not be read as deployment performance for the shipped head.
    \item Authenticity outputs: no concerns found, potentially synthetic, potentially edited, reused, inconsistent, or inconclusive.
    \item Real data: D-Fire originals. Synthetic training generators: SD 1.5 and SDXL; PixArt-alpha and Kolors held out.
    \item Evaluation normalization: longest side 512 px, JPEG quality 90; a validation/calibration threshold was tuned to 5\% real-image flags. This is not a production false-positive guarantee.
  \end{itemize}
  \vfill
  \scriptsize Advisory only: the experiment is narrow and does not establish detection of newer or closed generators, current incidents, location, or capture time.
\end{frame}

\begin{frame}{Image Evidence Assessment in TerraWatch}
  \centering
  \Large Authorized incident report
  \quad$\rightarrow$\quad Image upload
  \quad$\rightarrow$\quad Existing authenticity model
  \quad$\rightarrow$\quad Assessment result

  \vspace{1em}
  \Large$\downarrow$

  \vspace{1em}
  \Large Control Room / Range Officer review
  \quad$\rightarrow$\quad Human verification

  \vfill
  \small Automated assessment is stored separately from the incident status and never auto-rejects a report. The current field-note screen stores photo references locally; binary field-photo upload is not yet connected to that redesign flow.
\end{frame}

\begin{frame}{TerraWatch Dashboard: Role-Based User Experiences}
  \begin{columns}[T]
    \column{.25\textwidth}\textbf{COMMUNITY}\\[.4em]Public awareness, permitted community information, and responsible reporting.
    \column{.25\textwidth}\textbf{CITIZEN}\\[.4em]Public guidance and responsible reporting; no distinct citizen-account workflow is currently demonstrated.
    \column{.25\textwidth}\textbf{CONTROL ROOM}\\[.4em]Review incoming reports, inspect evidence, and coordinate authorized workflows.
    \column{.25\textwidth}\textbf{RANGE OFFICER}\\[.4em]Inspect authorized evidence and perform scoped review; the redesign field-note route is intended for field crew.
  \end{columns}
  \vfill
  \small Current repository status: staff incident reporting and field-note workflows exist; public/community reporting and binary field-photo persistence remain bounded by the existing authorization model.
\end{frame}

\begin{frame}{Integrated TerraWatch Architecture}
  \begin{description}
    \item[Ignition forecasting] Multimodal environmental data $\rightarrow$ TerraWatch Model $\rightarrow$ ranked ignition risk.
    \item[Fire-spread simulation] Ignition scenarios $\rightarrow$ Cell2Fire/Rothermel $\rightarrow$ saved spread examples.
    \item[Image authenticity] Report image $\rightarrow$ existing evidence assessor $\rightarrow$ advisory assessment.
  \end{description}
  \vfill
  \small These are separate processing paths. No evidence was found that the image assessor improves TWM or the spread simulator.
\end{frame}

\begin{frame}{Validation and Limitations}
  \begin{itemize}
    \item Synthetic-generator experiment: seen-generator and held-out-generator AUC are reported as 1.000 for clean inputs for both block-8 and final-layer features; these are clustered validation results over D-Fire originals and four stated diffusion families, not citizen-phone or current-incident evaluation.
    \item Real-source controls: final-layer real images flagged as synthetic were 1.6\% on FIgLib and 0.0\% on AoF; block-8 was less stable on FIgLib (59.2\%).
    \item Degradation matters: real-image flags rise under blur and half-resolution recompression; the 5\% threshold is only a calibration target.
    \item The fire/smoke head is not an authenticity verdict and can be fooled by generated fire images and burn scars.
    \item GCS contains research checkpoints and evaluation arrays, but no verified deployable Fire Evidence Assessor artifact is configured for TerraWatch. Run logs also leave OpenAI pretrained-weight provenance unresolved; runtime inference therefore remains unavailable until an approved artifact is configured.
  \end{itemize}
\end{frame}
\end{document}
